\subsection{Task 1-3: DCOPF Without Line Constraints}
% ----------------------------------------------------------

\subsubsection*{Optimal generation}

Dispatch follows the \emph{merit order} (cheapest generator first):

\begin{enumerate}
  \item Generator 1 is cheapest at 300~NOK/MWh (capacity 1000~MW).
        It covers the entire load of 900~MW on its own.
\end{enumerate}

\[
  P_{G1} = 900 \text{ MW}, \qquad P_{G2} = 0 \text{ MW}, \qquad P_{G3} = 0 \text{ MW}
\]

\subsubsection*{System (marginal) price}

The marginal generator is Generator 1, so
\[
  \lambda = c_1 = 300 \text{ NOK/MWh}
\]

\subsubsection*{Total system cost}

\[
  TC = 900 \times 300 = 270\,000 \text{ NOK}
\]

\subsubsection*{Phase angles}

Using $S_{\text{base}} = 100$~MVA, the net nodal injections in per unit are

\[
  P_1 = \frac{900-200}{100} = 7 \text{ p.u.}, \quad
  P_2 = \frac{0-200}{100}   = -2 \text{ p.u.}, \quad
  P_3 = \frac{0-500}{100}   = -5 \text{ p.u.}
\]

With $\delta_1 = 0$ as the reference, the nodal power balance
$P_n = \sum_l b_{nl}(\delta_n - \delta_l)$ gives two equations:

\textbf{Node 2} ($b_{12}=20,\; b_{23}=30$):
\[
  -2 = b_{12}(\delta_2 - 0) + b_{23}(\delta_2 - \delta_3)
     = 20\delta_2 + 30(\delta_2 - \delta_3)
  \;\Rightarrow\; 50\,\delta_2 - 30\,\delta_3 = -2
\]

\textbf{Node 3} ($b_{13}=10,\; b_{23}=30$):
\[
  -5 = b_{13}(\delta_3 - 0) + b_{32}(\delta_3 - \delta_2)
     = 10\delta_3 + 30(\delta_3 - \delta_2)
  \;\Rightarrow\; -30\,\delta_2 + 40\,\delta_3 = -5
\]

From the first equation:
\[
  \delta_2 = \frac{-2 + 30\,\delta_3}{50} = -0.04 + 0.6\,\delta_3
\]

Substituting into the second equation:
\[
  -30\!\left(-0.04 + 0.6\,\delta_3\right) + 40\,\delta_3 = -5
  \;\Rightarrow\; 1.2 - 18\,\delta_3 + 40\,\delta_3 = -5
  \;\Rightarrow\; 22\,\delta_3 = -6.2
\]

\[
  \delta_3 = -\frac{6.2}{22} = -\frac{31}{110} \approx -0.282 \text{ rad}
\]

\[
  \delta_2 = -0.04 + 0.6\times(-0.282) = -0.04 - 0.169 \approx -0.209 \text{ rad}
\]

\[
  \delta_1 = 0, \qquad \delta_2 \approx -0.209 \text{ rad}, \qquad \delta_3 \approx -0.282 \text{ rad}
\]

\subsubsection*{Power flows (in per unit, then convert)}

Power flows are defined as $P_{ij} = b_{ij}(\delta_i - \delta_j)$, where a positive
value means flow from bus $i$ to bus $j$.

\[
  P_{12} = b_{12}(\delta_1 - \delta_2) = 20\bigl(0 - (-0.209)\bigr)
         = 20 \times 0.209 \approx +4.18 \text{ p.u.}
         \;\Rightarrow\; +418 \text{ MW (bus 1 $\to$ bus 2)}
\]

\[
  P_{13} = b_{13}(\delta_1 - \delta_3) = 10\bigl(0 - (-0.282)\bigr)
         = 10 \times 0.282 \approx +2.82 \text{ p.u.}
         \;\Rightarrow\; +282 \text{ MW (bus 1 $\to$ bus 3)}
\]

\[
  P_{23} = b_{23}(\delta_2 - \delta_3) = 30\bigl(-0.209 - (-0.282)\bigr)
         = 30 \times 0.073 \approx +2.18 \text{ p.u.}
         \;\Rightarrow\; +218 \text{ MW (bus 2 $\to$ bus 3)}
\]

\noindent\textbf{Nodal balance check:}
\begin{align*}
  \text{Bus 1:} &\quad P_{12} + P_{13} = 4.18 + 2.82 = 7.00 \text{ p.u.} = P_1 \checkmark \\
  \text{Bus 2:} &\quad {-P_{12}} + P_{23} = -4.18 + 2.18 = -2.00 \text{ p.u.} = P_2 \checkmark \\
  \text{Bus 3:} &\quad {-P_{13}} - P_{23} = -2.82 - 2.18 = -5.00 \text{ p.u.} = P_3 \checkmark
\end{align*}

This is physically consistent: Generator~1 at bus~1 produces all the power, which
flows outward to the load buses. Bus~2 receives 418~MW, consumes 200~MW locally,
and forwards the remaining 218~MW onward to bus~3.

\subsubsection*{Nodal prices}

Since there are no binding transmission constraints, all nodes see the same
marginal price:
\[
  \lambda_1 = \lambda_2 = \lambda_3 = 300 \text{ NOK/MWh}
\]
